\documentclass[final,5p,times,twocolumn]{elsarticle}

\usepackage{amsmath,amssymb,amsfonts}
\usepackage{graphicx}
\usepackage{booktabs}
\usepackage{multirow}
\usepackage{algorithm}
\usepackage{algorithmic}
\usepackage{hyperref}
\usepackage{url}
\usepackage{color}
\usepackage{microtype}

\journal{Advances in Space Research}

\begin{document}

\begin{frontmatter}

\title{GMFRM-AreaNet: Global-Local Multiscale Feature Refinement Mamba Network for Remote Sensing Change Detection and Physical Area Quantification}

\author[inst1]{Author One\corref{cor1}}
\ead{author1@institution.edu}
\author[inst1]{Author Two}
\author[inst2]{Author Three}

\cortext[cor1]{Corresponding author.}
\affiliation[inst1]{organization={School of Geomatics and Remote Sensing},
            city={City},
            postcode={000000},
            state={State},
            country={Country}}
\affiliation[inst2]{organization={Department of Earth Observation and Geoinformatics},
            city={City},
            postcode={000000},
            state={State},
            country={Country}}

\begin{abstract}
Accurate change detection from bi-temporal remote sensing imagery is essential for monitoring urban expansion, agricultural shifts, and environmental changes. While deep learning models such as Convolutional Neural Networks (CNNs) and Vision Transformers (ViTs) are widely used, CNNs suffer from restricted receptive fields, and Transformers incur high quadratic computational costs on large satellite scenes. Recently, State-Space Models (SSMs), particularly Visual Mamba (VMamba), have shown promise by offering linear computational complexity with global receptive fields. However, existing Mamba-based methods often overlook fine boundary features in irregular parcels and provide only binary pixel masks without direct ground-area measurement. In this paper, we present GMFRM-AreaNet, an end-to-end framework combining a Global-Local Multiscale Feature Refinement Mamba Network with an automated physical area quantification module. Our model uses a weight-shared Siamese VMamba encoder to project bi-temporal images into a common feature space. At each feature stage, a Global-Local Context State-Space Block (GLCSB) couples global relational attention with local SSM convolutions to preserve fine textures and reduce false detections. A Multiscale Channel-Spatial Refinement Block (MCSRB) decoder aggregates contextual features across multiple receptive fields ($1\times1, 3\times3, 5\times5, 7\times7$) with parallel channel and spatial attention. We also implement an area calculation pipeline that directly converts predicted change masks into ground area measurements (square meters, hectares, and square kilometers) using the sensor Ground Sampling Distance (GSD) alongside transition matrices. We validate the framework on the high-resolution HRSCD benchmark ($0.5\text{ m}$ GSD) with 10,000 extracted patches, demonstrating reliable boundary delineation and accurate area statistics.
\end{abstract}

\begin{keyword}
Remote sensing \sep Change detection \sep Visual Mamba \sep State-Space Models \sep Multiscale refinement \sep Area quantification
\end{keyword}

\end{frontmatter}

\section{Introduction}
\label{sec:intro}
Remote sensing change detection identifies land surface modifications across the same geographical region observed at distinct timestamps $T_1$ and $T_2$. Tracking these alterations supports environmental governance, city planning, agricultural tracking, and disaster response \citep{ding2025survey}. Traditional change detection relied on manual thresholding, image differencing, and change vector analysis. While straightforward, these approaches are sensitive to radiometric discrepancies, shadows, seasonal illumination changes, and co-registration errors.

Deep learning improved feature representation through convolutional architectures \citep{daudt2018fully, fang2022snunet}. Convolutional networks extract local hierarchies effectively, but standard convolution kernels cannot easily model long-range spatial correlations across large satellite scenes. Vision Transformers (ViTs) addressed this by using global self-attention \citep{chen2022remote, bandara2022transformer}. However, Transformer attention scales quadratically ($\mathcal{O}(N^2)$) with token sequence length, which creates substantial memory and compute demands when handling high-resolution imagery.

State-Space Models (SSMs), such as Mamba and Visual Mamba (VMamba) \citep{chen2024changemamba}, offer an attractive middle ground: they deliver linear complexity ($\mathcal{O}(N)$) while maintaining full-image context. Even so, two practical gaps remain:
\begin{enumerate}
    \item \textbf{Boundary Smoothing and Detail Loss:} Long-sequence scans in Mamba can over-smooth boundaries of small, slender, or fragmented change parcels, leading to false negatives at parcel edges.
    \item \textbf{Absence of Direct Physical Area Metrics:} Standard change detection algorithms output binary mask grids ($0$ or $1$). For agricultural administrators and planners, raw pixel grids need to be converted into physical ground units (hectares, square kilometers) and land-cover transition matrices.
\end{enumerate}

To resolve these issues, we introduce \textbf{GMFRM-AreaNet}. Our primary contributions are:
\begin{itemize}
    \item A Siamese VMamba encoder coupled with a Global-Local Context State-Space Block (GLCSB) that balances global spatial context with local edge detail.
    \item A Multiscale Channel-Spatial Refinement Block (MCSRB) decoder using four parallel kernel sizes ($1\times1, 3\times3, 5\times5, 7\times7$) with dual channel and spatial attention to reconstruct clean parcel boundaries.
    \item An automated post-classification area quantification pipeline that uses sensor Ground Sampling Distance (GSD) and morphological filtering to report ground changes in physical units ($\text{m}^2$, $\text{ha}$, $\text{km}^2$) and transition dynamics.
    \item A complete reproducible Colab/PyTorch pipeline trained on 10,000 balanced patches from the French IGN HRSCD dataset, incorporating persistent Google Drive checkpoint caching for fault-tolerant training.
\end{itemize}

\section{Methodology}
\label{sec:method}

\subsection{Architecture Overview}
The overall structure of GMFRM-AreaNet is shown in Fig.~\ref{fig:arch}. Given bi-temporal images $I^{T_1}, I^{T_2} \in \mathbb{R}^{3 \times H \times W}$, a weight-shared Siamese VMamba encoder generates hierarchical feature maps $\{f^{T_1}_n, f^{T_2}_n\}_{n=1}^4$ at resolutions of $1/4, 1/8, 1/16$, and $1/32$. At each scale, bi-temporal features enter a Global-Local Context State-Space Block (GLCSB) to compute change representations. The decoder then aggregates these features through four Multiscale Channel-Spatial Refinement Blocks (MCSRB) to generate the final change probability map $\hat{Y} \in [0, 1]^{H \times W}$.

\begin{figure*}[!t]
\centering
% \includegraphics[width=0.92\linewidth]{figures/architecture.pdf}
\caption{System architecture of GMFRM-AreaNet, showing the Siamese VMamba encoder, multiscale GLCSB feature interaction modules, MCSRB decoder blocks, and the ground-area quantification pipeline.}
\label{fig:arch}
\end{figure*}

\subsection{Siamese Visual Mamba Encoder}
The encoder consists of four stages built from Visual State Space (VSS) blocks with 2D Selective Scan (2DSS). For an input feature $U$, 2DSS scans spatial tokens along four directional paths $d \in \{1, 2, 3, 4\}$:
\begin{itemize}
    \item $d=1$: Top-left to bottom-right (standard raster scan).
    \item $d=2$: Bottom-right to top-left (reversed raster scan).
    \item $d=3$: Top-right to bottom-left (column-wise scan).
    \item $d=4$: Bottom-left to top-right (reversed column-wise scan).
\end{itemize}
Each 1D sequence is processed by discrete state-space parameters $(\mathbf{A}, \mathbf{B}, \mathbf{V})$:
\begin{align}
    h_t &= \mathbf{A} h_{t-1} + \mathbf{B} x_t \\
    y_t &= \mathbf{V} h_t
\end{align}
The four directional outputs are reshaped and fused:
\begin{equation}
    T = \text{Merge}(y_1, y_2, y_3, y_4)
\end{equation}
This allows each spatial token to interact with all surrounding pixels with linear computational complexity.

\subsection{Global-Local Context State-Space Block (GLCSB)}
To capture changes between $T_1$ and $T_2$, we concatenate bi-temporal features along channels: $D = [f^{T_1}_n, f^{T_2}_n] \in \mathbb{R}^{2C \times H_n \times W_n}$. GLCSB splits processing into two complementary paths:
\begin{enumerate}
    \item \textbf{Global Relational Path:} Projects $D$ into Query ($Q$), Key ($K$), and Value ($V$) representations. Spatial correlation weights $W$ and the global feature $F_{\text{Global}}$ are computed as:
    \begin{align}
        W &= \text{Softmax}\left(\frac{QK^\top}{\sqrt{d_k}}\right) \\
        F_{\text{Global}} &= W \cdot V + \phi_d(D)
    \end{align}
    where $\phi_d$ is a $1\times1$ residual projection.
    \item \textbf{Local Enhancement Path:} Runs $D$ through a 2DSS layer followed by a $3\times3$ depth-aware convolution and ReLU activation to preserve edge details:
    \begin{equation}
        F_{\text{Local}} = \text{ReLU}\left(\text{Conv}_{3\times3}(\text{SSM}(D))\right)
    \end{equation}
\end{enumerate}
The outputs of both paths are combined with a $1\times1$ convolution:
\begin{equation}
    F = \text{Conv}_{1\times1}\left([F_{\text{Global}}, F_{\text{Local}}]\right)
\end{equation}

\subsection{Multiscale Channel-Spatial Refinement Block (MCSRB)}
In the decoder, feature maps pass through four parallel convolution branches with kernel sizes $1\times1, 3\times3, 5\times5$, and $7\times7$:
\begin{align}
    S_1 &= \text{Conv}_{1\times1}(F), \quad S_2 = \text{Conv}_{3\times3}(F) \\
    S_3 &= \text{Conv}_{5\times5}(F), \quad S_4 = \text{Conv}_{7\times7}(F) \\
    J &= [S_1, S_2, S_3, S_4] \in \mathbb{R}^{4C \times H \times W}
\end{align}
To refine this concatenated representation:
\begin{itemize}
    \item \textbf{Channel Attention:} Global average pooling extracts channel descriptors $z = \text{GAP}(J)$, followed by a two-layer bottleneck MLP:
    \begin{equation}
        O_C = \sigma\left(\text{Conv}_{1\times1}(\text{ReLU}(\text{Conv}_{1\times1}(z)))\right), \quad J_C = O_C \otimes J
    \end{equation}
    \item \textbf{Spatial Attention:} Max-pooling and average-pooling across channels are concatenated and passed through a $1\times1$ convolution:
    \begin{equation}
        O_S = \sigma\left(\text{Conv}_{1\times1}([J_{\max}, J_{\text{avg}}])\right), \quad J_S = O_S \otimes J
    \end{equation}
\end{itemize}
The refined block output is:
\begin{equation}
    X = \text{ReLU}\left(\text{BN}\left(\text{Conv}_{1\times1}(J + J_C + J_S)\right)\right)
\end{equation}

\subsection{Adaptive Objective Loss}
Satellite change detection datasets typically have far more unchanged pixels than changed pixels. To address this class imbalance without manual weight tuning, we use an adaptive combination of Binary Cross-Entropy ($L_{\text{BCE}}$) and Dice loss ($L_{\text{Dice}}$):
\begin{equation}
    L_{\text{Total}} = \alpha L_{\text{BCE}} + \beta L_{\text{Dice}}
\end{equation}
The weighting factors $\alpha$ and $\beta$ dynamically adjust based on loss magnitude:
\begin{equation}
    \alpha = \frac{L_{\text{BCE}}}{L_{\text{BCE}} + L_{\text{Dice}} + \epsilon}, \quad \beta = \frac{L_{\text{Dice}}}{L_{\text{BCE}} + L_{\text{Dice}} + \epsilon}
\end{equation}
where $\epsilon = 10^{-5}$ prevents division by zero. This gives higher relative weight to the loss term that is currently harder to optimize during training.

\subsection{Ground Area Quantification Framework}
\label{sec:area_calc}
To convert binary predictions $\hat{Y}_{i,j} \in \{0, 1\}$ into physical units, we use the Ground Sampling Distance (GSD) $\Delta x$ and $\Delta y$ in meters. For HRSCD orthophotos, $\Delta x = \Delta y = 0.5\text{ m}$, so each pixel represents:
\begin{equation}
    A_{\text{pixel}} = \Delta x \cdot \Delta y = 0.25\text{ m}^2
\end{equation}

\subsubsection{Physical Area Estimation}
The total changed area is obtained by summing all changed pixels:
\begin{equation}
    \mathcal{A}_{\text{change}} = \left(\sum_{i=1}^{H}\sum_{j=1}^{W} \hat{Y}_{i,j}\right) \times A_{\text{pixel}} \quad [\text{m}^2]
\end{equation}
In hectares ($\text{ha}$) and square kilometers ($\text{km}^2$):
\begin{align}
    \mathcal{A}_{\text{ha}} &= \frac{\mathcal{A}_{\text{change}}}{10{,}000} = 10^{-4} \sum_{i,j} \hat{Y}_{i,j} A_{\text{pixel}} \\
    \mathcal{A}_{\text{km}^2} &= \frac{\mathcal{A}_{\text{change}}}{1{,}000{,}000} = 10^{-6} \sum_{i,j} \hat{Y}_{i,j} A_{\text{pixel}}
\end{align}

\subsubsection{Morphological Filtering and Parcel Grouping}
Raw neural network predictions may contain scattered single-pixel noise. We apply a $3\times3$ morphological opening to eliminate isolated false alarms:
\begin{equation}
    \hat{Y}_{\text{clean}} = (\hat{Y} \circ K) = (\hat{Y} \ominus K) \oplus K
\end{equation}
Connected component analysis then segments the binary mask into discrete geographic parcels, computing parcel count, individual parcel areas, perimeters, and centroids.

\subsubsection{Multiclass Transition Matrix}
When land-cover maps $C^{T_1}, C^{T_2} \in \{1, \dots, K\}$ are available across $K$ classes (e.g., Urban, Cropland, Forest, Wetland, Water), we compute a $K \times K$ transition matrix $\mathbf{M}$ in hectares:
\begin{equation}
    \mathbf{M}_{u,v} = \frac{A_{\text{pixel}}}{10{,}000} \sum_{i=1}^H \sum_{j=1}^W \mathbb{I}\left(C^{T_1}_{i,j} = u \land C^{T_2}_{i,j} = v \land \hat{Y}_{i,j} = 1\right)
\end{equation}
This records the exact physical area converted from class $u$ to class $v$.

\section{Experimental Setup}
\label{sec:experiments}

\subsection{Dataset and Preprocessing}
We evaluate GMFRM-AreaNet on the French National Institute of Geographic and Forest Information (IGN) HRSCD benchmark. The dataset contains aerial orthophotos acquired in 2006 ($T_1$) and 2012 ($T_2$) at $0.5\text{ m}$ GSD with pixel-level change masks. 

We matched 291 large orthophoto tile sets ($10{,}000 \times 10{,}000$ pixels) across $T_1$, $T_2$, and change labels. Using a sliding window of $256 \times 256$ pixels, we extracted a curated dataset of 10,000 paired patches. Pure background patches with no changes were subsampled to maintain a balanced ratio between changed and unchanged regions. The 10,000 patches were split into 8,500 training patches ($85\%$) and 1,500 validation patches ($15\%$), and the split was indexed into a persistent JSON cache to ensure exact reproducibility across runs.

\subsection{Implementation and Training Details}
The model was implemented in PyTorch and trained on an NVIDIA T4 GPU (16 GB VRAM) in Google Colab. Training parameters were set as follows:
\begin{itemize}
    \item Input image dimensions: $256 \times 256 \times 3$.
    \item Optimizer: Adam with initial learning rate $\eta = 1 \times 10^{-4}$ and weight decay $\lambda = 5 \times 10^{-3}$.
    \item Scheduler: Cosine Annealing learning rate schedule over 60 epochs down to $\eta_{\min} = 1 \times 10^{-6}$.
    \item Batch size: 16 with automatic mixed precision (AMP FP16).
    \item Data augmentation: Random horizontal and vertical flips, random $90^\circ$ rotations, and ImageNet channel normalization.
    \item Checkpointing: The training pipeline saves a complete checkpoint (`checkpoint\_latest.pth`) after every epoch directly to Google Drive. If a runtime disconnection occurs, the training script automatically restores the epoch, model weights, optimizer, and scheduler states to continue without loss of progress.
\end{itemize}

\subsection{Evaluation Metrics}
We measure performance using standard change detection metrics: Overall Accuracy (OA), Intersection over Union (IoU), Precision (Pre), Recall (Rec), and F1-score:
\begin{align}
    \text{OA} &= \frac{\text{TP} + \text{TN}}{\text{TP} + \text{FP} + \text{FN} + \text{TN}} \\
    \text{Pre} &= \frac{\text{TP}}{\text{TP} + \text{FP}}, \quad \text{Rec} = \frac{\text{TP}}{\text{TP} + \text{FN}} \\
    \text{F1} &= \frac{2 \times \text{Pre} \times \text{Rec}}{\text{Pre} + \text{Rec}}, \quad \text{IoU} = \frac{\text{TP}}{\text{TP} + \text{FP} + \text{FN}}
\end{align}

\subsection{Quantitative Results}
Table~\ref{tab:comparison} compares GMFRM-AreaNet against standard CNN, Transformer, and Mamba baselines on high-resolution change detection benchmarks.

\begin{table*}[!t]
\centering
\caption{Quantitative comparison of change detection models on high-resolution benchmarks. Bold indicates the best result.}
\label{tab:comparison}
\begin{tabular}{lcccccc}
\toprule
Model & Architecture & Params (M) & FLOPs (G) & OA (\%) & IoU (\%) & F1 (\%) \\
\midrule
FC-EF \citep{daudt2018fully} & CNN & 1.35 & 3.58 & 94.57 & 46.74 & 63.71 \\
SNUNet \citep{fang2022snunet} & DenseNet & 50.44 & 54.83 & 93.88 & 39.95 & 57.09 \\
BIT \citep{chen2022remote} & ResNet+ViT & 17.99 & 15.75 & 95.23 & 51.44 & 67.93 \\
ChangeFormer \citep{bandara2022transformer} & Pure ViT & 29.75 & 21.18 & 94.86 & 50.25 & 66.89 \\
ChangeMamba \citep{chen2024changemamba} & VMamba & 42.93 & 25.75 & 94.90 & 48.31 & 65.15 \\
CDMamba \citep{zhang2025cdmamba} & Mamba+CNN & 21.87 & 59.83 & 94.51 & 45.97 & 62.98 \\
\textbf{GMFRM-AreaNet (Ours)} & \textbf{Siamese VMamba} & \textbf{26.15} & \textbf{30.46} & \textbf{97.43} & \textbf{68.99} & \textbf{81.65} \\
\bottomrule
\end{tabular}
\end{table*}

\subsection{Ablation Study}
To assess the contribution of the proposed modules, we conducted an ablation study summarized in Table~\ref{tab:ablation}. Incorporating GLCSB improves contextual discrimination between true changes and seasonal differences. Adding MCSRB sharpens boundary definition for narrow parcels, and combining both blocks yields the highest overall accuracy and IoU.

\begin{table}[!h]
\centering
\caption{Ablation analysis of individual architectural components.}
\label{tab:ablation}
\begin{tabular}{lcccc}
\toprule
Configuration & GLCSB & MCSRB & IoU (\%) & F1 (\%) \\
\midrule
Baseline Backbone & -- & -- & 67.29 & 80.44 \\
Base + GLCSB & \checkmark & -- & 67.66 & 80.70 \\
Base + MCSRB & -- & \checkmark & 68.81 & 81.52 \\
\textbf{GMFRM-AreaNet (Full)} & \checkmark & \checkmark & \textbf{68.99} & \textbf{81.65} \\
\bottomrule
\end{tabular}
\end{table}

\section{Discussion and Practical Considerations}
\label{sec:discussion}
Training VMamba models on consumer-grade hardware involves specific memory tradeoffs. On an NVIDIA T4 GPU (16 GB), batch sizes of 16 run efficiently with automatic mixed precision, completing each training epoch across 10,000 patches in approximately 50 to 60 seconds. On laptop GPUs with 4 GB VRAM (such as the RTX 3050), smaller batch sizes (batch size 4) paired with gradient accumulation are required to avoid out-of-memory errors. The persistent checkpoint caching strategy implemented here protects against cloud session timeouts, making practical training feasible even on free cloud tiers.

\section{Conclusion}
\label{sec:conclusion}
In this paper, we introduced GMFRM-AreaNet for high-resolution remote sensing change detection. The network pairs a weight-shared Siamese Visual Mamba encoder with GLCSB feature interaction modules and an MCSRB multi-kernel decoder. This design captures broad spatial context while preserving parcel boundary details. An integrated area quantification module converts pixel predictions directly into ground-area measurements and transition matrices. Experiments on the HRSCD benchmark confirm that the framework provides both strong detection performance and actionable spatial metrics for land management.

\bibliographystyle{elsarticle-num}
\begin{thebibliography}{10}

\bibitem{ding2025survey}
L.~Ding, D.~Hong, M.~Zhao, H.~Chen, et~al.,
\newblock A survey of sample-efficient deep learning for change detection in remote sensing: Tasks, strategies, and challenges,
\newblock \emph{IEEE Geoscience and Remote Sensing Magazine}, vol.~13, pp.~164--189, 2025.

\bibitem{daudt2018fully}
R.~C. Daudt, B.~Le~Saux, and A.~Boulch,
\newblock Fully convolutional siamese networks for change detection,
\newblock in \emph{IEEE International Conference on Image Processing (ICIP)}, 2018, pp.~4063--4067.

\bibitem{fang2022snunet}
S.~Fang, K.~Li, J.~Shao, and Z.~Li,
\newblock SNUNet-CD: A densely connected siamese network for change detection of VHR images,
\newblock \emph{IEEE Geoscience and Remote Sensing Letters}, vol.~19, pp.~1--5, 2022.

\bibitem{chen2022remote}
H.~Chen, Z.~Qi, and Z.~Shi,
\newblock Remote sensing image change detection with transformers,
\newblock \emph{IEEE Transactions on Geoscience and Remote Sensing}, vol.~60, pp.~1--14, 2022.

\bibitem{bandara2022transformer}
W.~G.~C. Bandara and V.~M. Patel,
\newblock A transformer-based siamese network for change detection,
\newblock in \emph{IEEE International Geoscience and Remote Sensing Symposium (IGARSS)}, 2022, pp.~207--210.

\bibitem{chen2024changemamba}
H.~Chen, J.~Song, C.~Han, J.~Xia, and N.~Yokoya,
\newblock ChangeMamba: Remote sensing change detection with spatiotemporal state space model,
\newblock \emph{IEEE Transactions on Geoscience and Remote Sensing}, vol.~62, 2024.

\bibitem{zhang2025cdmamba}
H.~Zhang, K.~Chen, C.~Liu, H.~Chen, Z.~Zou, and Z.~Shi,
\newblock CDMamba: Incorporating local clues into Mamba for remote sensing image binary change detection,
\newblock \emph{IEEE Transactions on Geoscience and Remote Sensing}, vol.~63, 2025.

\end{thebibliography}

\end{document}
